
%\needspace{4\baselineskip}
\color{uniblue}{\subsection{Дифференциальная токовая защита трансформатора (ДЗТ)}}
\color{black}

\begin{enumerate}[label=\arabic{section}.\arabic{subsection}.\arabic{enumi}, labelsep=4pt, leftmargin=0pt, itemindent=57pt, itemsep=0pt, parsep=5pt] % если что можно крутить itemsep
%%%%%%%%%%%%%%%%%%%%%%%%%%%%%%%%%%%%%%%%%%%%%%%%%%%%%%%%%%%%%%%%%%%%%%  ОБЩИЕ СВЕДЕНИЯ %%%%%%%%%%%%%%%%%%%%%%%%%%%%%%%%%%%%%%%%%%%%%%%%%%
\item Общие сведения

\begin{enumerate}[label=\arabic{section}.\arabic{subsection}.\arabic{enumi}.\arabic*, labelsep=4pt, leftmargin=0em, itemindent=65pt, parsep=0pt]

\item
Дифференциальная токовая защита применяется в качестве основной защиты двухобмоточного или трехобмоточного силового трансформатора. ДЗТ обладает абсолютной селективностью и высоким быстродействием. Работа дифференциальных токовых функций не зависит от режима заземления нейтрали сети, и, следовательно, ДЗТ  может быть использована в качестве основной защиты от всех видов коротких замыканий.

\item
В устройстве дифференциальная защита трансформатора представлена функциональным блоком «ДЗТ». Функциональный блок содержит две независимые друг от друга защитные функции:
\begin{itemize}
\item дифференциальный токовый орган с торможением (ДТЗт);
\item дифференциальная токовая отсечка (ДТО).
\end{itemize}
\item
Функциональный блок «ДЗТ» также включает в себя вспомогательные функции, необходимые для правильной работы в различных режимах:
\begin{itemize}
    \item детектор второй гармоники (Д2Г);
    \item детектор пятой гармоники (Д5Г);
    \item функция блокировки при внешних коротких замыканиях (БВКЗ);
    \item функция контроля токовых цепей по величине небаланса (КЦТнеб).
\end{itemize}

\item
Основные характеристики ДЗТ:
\begin{itemize}
    \item полное время срабатывания чувствительной ступени ДЗТ при дифференциальном токе более двухкратного тока уставки, в том числе с учетом наличия апериодической составляющей, с учетом времени срабатывания выходного реле - не более 50 мс.
\end{itemize}

\item
Оперативный вывод ДЗТ из работы происходит при появлении активного сигнала на входе «ОВ ДЗТ», при этом осуществляется вывод защитных функций ДТО и ДТЗт. 
\end{enumerate}
%%%%%%%%%%%%%%%%%%%%%%%%%%%%%%%%%%%%%%%%%%%%%%%%%%%%%%%%%%%%%%%%%%%%%%  ПРИНЦИП ДЕЙСТВИЯ %%%%%%%%%%%%%%%%%%%%%%%%%%%%%%%%%%%%%%%%%%%%%%%%%%
%\color{uniblue}{\subsubsection{Принцип действия}}
%\color{black}
\item
Принцип работы дифференциальной защиты основан на первом правиле Кирхгофа. Защищаемый объект рассматривается как электрический узел, а его связи с электрической сетью как ветви узла.

%%%%%%%%%%%%%%%%%%%%%%%%%%%%%%%%%%%%%%%%%%%%%%%%%%%%%%%%%%%%%%%%%%%%%% ЦИФРОВОЕ ВЫРАВНИВАНИЕ %%%%%%%%%%%%%%%%%%%%%%%%%%%%%%%%%%%%%%%%%%%%%%%%%%
%\item Цифровое выравнивание подведенных к устройству токов, исключение утроенного тока нулевой последовательности\\
%\color{uniblue}{\subsubsection{Цифровое выравнивание подведенных к устройству токов, исключение утроенного тока нулевой последовательности}}
%\color{black}

\item  Цифровое выравнивание, исключение токов нулевой последовательности

\begin{enumerate}[label=\arabic{section}.\arabic{subsection}.\arabic{enumi}.\arabic*, labelsep=4pt, leftmargin=0em, itemindent=65pt, parsep=0pt]

\item
Силовые трансформаторы имеют различные группы соединений обмоток, при этом у измерительных ТТ, установленных по сторонам защищаемого трансформатора, различные номинальные значения первичных и вторичных токов, поэтому, прежде чем выполнить сравнение фазных токов каждой из сторон трансформатора, необходимо выполнить выравнивание измеряемых токов по модулю (амплитуде), по фазе и, при необходимости, выполнить фильтрацию токов нулевой последовательности.
Устройство по заданным параметрам осуществляет выравнивание измеренных токов сторон в соответствии с формулой:

% ФОРМУЛА №1
\begin{equation*}
    \begin{bmatrix}
        \vec I_{A,n*} \\ 
        \vec I_{B,n*} \\ 
        \vec I_{C,n*}
    \end{bmatrix} 
    = 
    k_{\textsl{П},n} \cdot k_{AM,n} \cdot 
    \begin{bmatrix}
        M_{11,n} & M_{12,n} & M_{13,n} \\ 
        M_{21,n} & M_{22,n} & M_{23,n} \\ 
        M_{31,n} & M_{32,n} & M_{33,n}
    \end{bmatrix} 
    \cdot 
    \begin{bmatrix}
        \vec I'_{A,n} \\ 
        \vec I'_{B,n} \\ 
        \vec I'_{C,n}
    \end{bmatrix},
\end{equation*}

\begin{eqexpl}[25mm]
    \item{$n$} индекс, обозначающий сторону ДЗТ;
    \item{$k_{\textsl{П},n}$} коэффициент коррекции полярности;
    \item{$k_{AM,n}$} коэффициент амплитудного выравнивания;
    \item{$M_{11,n}\dots M_{33,n}$} коэффициенты матрицы компенсации фазового сдвига в соответствии с группой соединения обмоток трансформатора;
    \item{$\vec I'_{A,n},\vec I'_{B,n},\vec I'_{C,n}$} фазные токи, приведённые к номиналу токового входа терминала, в о.е. Вычисляются по формулам:
    \begin{align*}
        \vec I'_{A,n} &= \frac{\vec I_{A,n}}{I_{\textsl{ном.терм.n}}}, &
        \vec I'_{B,n} &= \frac{\vec I_{B,n}}{I_{\textsl{ном.терм.n}}}, &
        \vec I'_{C,n} &= \frac{\vec I_{C,n}}{I_{\textsl{ном.терм.n}}},
    \end{align*}
    где {$\vec I_{A,n},\vec I_{B,n},\vec I_{C,n}$} -- измеренные фазные токи стороны n (комплексные, без учета коэффициента полярности), А.\\
    {$I_{\textsl{ном.терм.n}}$} --  номинальный ток токового входа терминала выбранного ответвления, подключенного к трансформатору току стороны n, А.

\end{eqexpl}

\item
Коэффициент коррекции полярности ($k_{\textsl{П},n}$) предназначен для учета различных схем сборки «звезды» трансформаторов тока по плечам дифференциальной токовой защиты трансформатора для правильного выравнивания токов. Если обмотки ТТ соединены противофазно, то коэффициент принимается равным 1, иначе (при синфазном подключении), коэффициент принимается равным -1 (см. рис. \ref{dzt:connection}). 

\vspace{1mm}
\begin{figure}[H]
\centering
\includegraphics[width=0.5\textwidth,height=0.5\textheight,keepaspectratio]{\fbpath/17. ДЗТ 35/_latex/img/connection.png}
\captionof{figure}{Поясняющая схема включения трансформаторов тока: (а) противофазное включение ТТ (б) синфазное включение ТТ}\label{dzt:connection}
\end{figure}

\item
Расчет коэффициента амплитудного выравнивания ($k_{AM,n}$) производится по задаваемым параметрам величин номинальных значений силового трансформатора и измерительных трансформаторов тока в соответствии с выражениями:

% ФОРМУЛА №2
\begin{equation*}
k_{AM,n}= \frac{I_{\textsl{перв,n}}\cdot  I_{\textsl{ном.терм,n}} }{I_{\textsl{баз,n}}\cdot I_{\textsl{втор,n}}\cdot k_{\textsl{сх,n}}},
\end{equation*}
\begin{equation*}
I_{\textsl{баз,n}}=\frac{S_{\textsl{баз}}}{\sqrt{3} \cdot U_{\textsl{баз,n}} },
\end{equation*}
\begin{eqexpl}[25mm]
\item{$n$} индекс, обозначающий сторону (плечо) ДЗТ;
\item{$I_{\textsl{перв,n}}$}номинальный первичный ток ТТ стороны $n$;
\item{$I_{\textsl{втор,n}}$}номинальный вторичный ток ТТ стороны $n$;
\item{$I_{\textsl{ном.терм,n}}$}номинальное значение тока аналоговых входов устройства, подключенных к ТТ стороны $n$;
\item{$k_{\textsl{сх,n}}$}коэффициент, учитывающий схему соединения обмоток ТТ стороны~$n$. Данный параметр принимается равным 1, когда ТТ стороны $n$ собраны в звезду, или равным $\sqrt{3}$, когда ТТ стороны $n$ собраны в треугольник. Коэффициенты схемы соединения ТТ для каждой из сторон задается уставками «KсхемСт1», «KсхемСт2», «KсхемСт3»;
\item{$S_{\textsl{баз}}$}базисная мощность, равная номинальной мощности силового трансформатора;
\item{$U_{\textsl{баз,n}}$}базисное напряжение, равное номинальному напряжению обмотки трансформатора соответствующей стороне $n$.
\end{eqexpl}

%\vspace{1mm}
%\emph{\textls[50]{Выравнивание токов сторон по фазе}}
%\par
%\vspace{1mm}
\item 
Для компенсации фазового сдвига, обусловленного различными группами соединения обмоток трансформатора, для каждой из сторон предусмотрены задаваемые параметры («Nсх1», «Nсх2», «Nсх3»). Значения данных параметров (от 0 до 11) определяют коэффициенты матрицы $M_{11,n}\dots M_{33,n}$ в соответствие с таблицей \ref{dzt:tbl0}.

%\vspace{1mm}
%\emph{\textls[50]{Компенсация утроенного тока нулевой последовательности}}
%\par
%\vspace{1mm}

\item 
Для обмоток силового трансформатора, соединенных в звезду с заземленной нейтралью, и, если не используется других мер исключения токов нулевой последовательности (соединение вторичных обмоток ТТ в треугольник) всегда рекомендуется предусматривать исключение составляющих нулевой последовательности из фазных токов. Компенсация токов нулевой последовательности сторон задается параметрами «Компенс3I0ст1», «Компенс3I0ст2», «Компенс3I0ст3». В случае, когда параметр компенсации токов нулевой последовательности для стороны равен «0» (без компенсации тока $3I0$), то в формуле для выравнивания измеряемых токов сторон применяются коэффициенты матрицы M (см. таблицу \ref{dzt:tbl0}). В случае, когда параметр компенсации токов нулевой последовательности для стороны равен «1» (с компенсацией токов $3I0$), то в формуле для выравнивания измеряемых токов сторон применяются коэффициенты матрицы M0 (см. таблицу \ref{dzt:tbl0}).

%%%%%%%%%%%%%%%%%%%%%%%%%%%%%%%%%%%%%%%%%%%%%%%%%%%%%%%%%%%%%%%%%%%%%%%%%% ТАБЛИЦА 1 %%%%%%%%%%%%%%%%%%%%%%%%%%%%%%%%%%%%%%%%%%%%%%%%%%%%%%%%%%%%%%%%%%%
%\setlength{\extrarowheight}{0.0cm} % высота ячеек в таблице
\small
\begin{longtable}{|>{\centering\arraybackslash}m{1.7cm}|>{\centering\arraybackslash}m{4cm}|>{\centering\arraybackslash}m{3cm}|>{\centering\arraybackslash}m{4cm}|>{\centering\arraybackslash}m{3cm}|}
\caption{Компенсация группы соединения \hfill\vspace{-0.5\baselineskip}}\label{dzt:tbl0}\\
\hhline{|-|-|-|-|-|}
\rowcolor{gray!30}
\multicolumn{1}{|c|}{} & \multicolumn{2}{c|}{Без компенсации токов 3I0} & \multicolumn{2}{c|}{С компенсацией токов 3I0} \\
\hhline{|-|-|-|-|-|}
\rowcolor{gray!30}
Векторная группа & Матрица М & Группы соединений обмоток & Матрица М0 & Группы соединений обмоток \\ 
\hhline{|-|-|-|-|-|}
\endfirsthead

%\caption{\emph{Продолжение\hfill\vspace{-0.5\baselineskip}}} \\ % Отступ обозначения таблицы от таблицы и выравнивание надписи слева
\caption*{\hspace{3pt}\emph{Продолжение таблицы \ref{dzt:tbl0}\hfill\vspace{-0.5\baselineskip}}} \\ % сделано по ГОСТ 2.105 п.6.8.7
\hhline{|-|-|-|-|-|}
\rowcolor{gray!30}
\multicolumn{1}{|c|}{} & \multicolumn{2}{c|}{Без компенсации токов 3I0} & \multicolumn{2}{c|}{С компенсацией токов 3I0} \\
\hhline{|-|-|-|-|-|}
\rowcolor{gray!30}
Векторная группа & Матрица М & Группы соединений обмоток & Матрица М0 & Группы соединений обмоток \\ 
\hhline{|-|-|-|-|-|}
\endhead
\hhline{|-|-|-|-|-|}
\endfoot
%\hline
\endlastfoot

\centering 0 & \centering $$\begin{vmatrix}1&0&0\\0&1&0\\0&0&1\end{vmatrix}$$ & \centering \parbox{1cm}{Yy0 Dd0} & \centering $$\displaystyle\frac{1}{3}\times \begin{vmatrix}2&-1&-1\\-1&2&-1\\-1&-1&2\end{vmatrix}$$ & \centering\arraybackslash \parbox{1cm}{Yyn0 YNy0} \\
\hline

\centering 1 & \centering $$\displaystyle\frac{1}{\sqrt{3}}\times\begin{vmatrix}2&-1&-1\\-1&2&-1\\-1&-1&2\end{vmatrix}$$ & \centering \parbox{1cm}{Yd1 Dy1} & \centering Равна матрице №1 & \centering\arraybackslash \parbox{1cm}{YNd1 Dyn1} \\
\hline

\centering 2 & \centering $$\begin{vmatrix}0&-1&0\\0&0&-1\\-1&0&0\end{vmatrix}$$ & \centering \parbox{1cm}{Yy2 Dd2} & \centering $$\displaystyle\frac{1}{3}\times \begin{vmatrix}1&-2&1\\1&1&-2\\-2&1&1\end{vmatrix}$$ & \centering\arraybackslash \parbox{1cm}{YNy2 Yyn2} \\
\hline

\centering 3 & \centering $$\displaystyle\frac{1}{\sqrt{3}}\times\begin{vmatrix}0&-1&1\\1&0&-1\\-1&1&0\end{vmatrix}$$ & \centering \parbox{1cm}{Yd3 Dy3} & \centering Равна матрице №3 & \centering\arraybackslash \parbox{1cm}{YNd3 Dyn3} \\
\hline

\centering 4 & \centering $$\begin{vmatrix}0&0&1\\1&0&0\\0&1&0\end{vmatrix}$$ & \centering \parbox{1cm}{Yy4 Dd4} & \centering $$\displaystyle\frac{1}{3}\times \begin{vmatrix}-1&-1&2\\2&-1&-1\\-1&2&-1\end{vmatrix}$$ & \centering\arraybackslash \parbox{1cm}{YNy4 Yyn4} \\
\hline

\centering 5 & \centering $$\displaystyle\frac{1}{\sqrt{3}}\times\begin{vmatrix}-1&0&1\\1&-1&0\\0&1&-1\end{vmatrix}$$ & \centering \parbox{1cm}{Yd5 Dy5} & \centering Равна матрице №5 & \centering\arraybackslash \parbox{1cm}{YNd5 Dyn5} \\
\hline

\centering 6 & \centering $$\begin{vmatrix}-1&0&0\\0&-1&0\\0&0&-1\end{vmatrix}$$ & \centering \parbox{1cm}{Yy6 Dd6} & \centering $$\displaystyle\frac{1}{3}\times \begin{vmatrix}-2&1&1\\1&-2&1\\1&1&-2\end{vmatrix}$$ & \centering\arraybackslash \parbox{1cm}{YNy6 Yyn6} \\
\hline

\centering 7 & \centering $$\displaystyle\frac{1}{\sqrt{3}}\times\begin{vmatrix}-1&1&0\\0&-1&1\\1&0&-1\end{vmatrix}$$ & \centering \parbox{1cm}{Yd7 Dy7} & \centering Равна матрице №7 & \centering\arraybackslash \parbox{1cm}{YNd7 Dyn7} \\
\hline

\centering 8 & \centering $$\begin{vmatrix}0&1&0\\0&0&1\\1&0&0\end{vmatrix}$$ & \centering \parbox{1cm}{Yy8 Dd8} & \centering $$\displaystyle\frac{1}{3}\times \begin{vmatrix}-1&2&-1\\-1&-1&2\\2&-1&-1\end{vmatrix}$$ & \centering\arraybackslash \parbox{1cm}{YNy8 Yyn8} \\
\hline

\centering 9 & \centering $$\displaystyle\frac{1}{\sqrt{3}}\times\begin{vmatrix}0&1&-1\\-1&0&1\\1&-1&0\end{vmatrix}$$ & \centering \parbox{1cm}{Yd9 Dy9} & \centering Равна матрице №9 & \centering\arraybackslash \parbox{1cm}{YNd9 Dyn9} \\
\hline

\centering 10 & \centering $$\begin{vmatrix}0&0&-1\\-1&0&0\\0&-1&0\end{vmatrix}$$ & \centering \parbox{1cm}{Yy10 Dd10} & \centering $$\displaystyle\frac{1}{3}\times \begin{vmatrix}1&1&-2\\-2&1&1\\1&-2&1\end{vmatrix}$$ & \centering\arraybackslash \parbox{1cm}{YNy10 Yyn10} \\
\hline

\centering 11 & \centering $$\displaystyle\frac{1}{\sqrt{3}}\times\begin{vmatrix}1&0&-1\\-1&1&0\\0&-1&1\end{vmatrix}$$ & \centering \parbox{1cm}{Yd11 Dy11} & \centering Равна матрице №11 & \centering\arraybackslash \parbox{1cm}{YNd11 Dyn11} \\
\hline

\end{longtable}
\normalsize

\end{enumerate}
%%%%%%%%%%%%%%%%%%%%%%%%%%%%%%%%%%%%%%%%%%%%%%%%%%%%%%%%%%%%%%%%%%%%%%  ДИФОРГАН С ТОРМОЖЕНИЕМ %%%%%%%%%%%%%%%%%%%%%%%%%%%%%%%%%%%%%%%%%%%%%%%%%%

%\needspace{3\baselineskip}
%\color{uniblue}{\subsubsection{Дифференциальный токовый орган с торможением (ДТЗт)}\label{sec:dtzt}}
%\color{black}

\item\label{sec:dtzt} Дифференциальный токовый орган с торможением (ДТЗт)

\begin{enumerate}[label=\arabic{section}.\arabic{subsection}.\arabic{enumi}.\arabic*, labelsep=4pt, leftmargin=0em, itemindent=65pt, parsep=0pt]

\item
Характеристика срабатывания ДТЗт в координатах $I_{\textsl{диф}}$ и $I_{\textsl{торм}}$ представляет собой ломаную линию (см. рисунок \ref{dzt:char}), состоящую из трех участков. Данная характеристика срабатывания обеспечивает чувствительность ДТЗт при малых токах КЗ в зоне действия защиты.

\vspace{1mm}
\begin{figure}[H]
\centering
\includegraphics[width=0.55\textwidth,height=0.55\textheight,keepaspectratio]{\fbpath/17. ДЗТ 35/_latex/img/chardtz.pdf}
\captionof{figure}{Характеристика срабатывания ДЗТ}\label{dzt:char}
\end{figure}

\item
Первый участок выполнен горизонтальным и определяется уставкой минимального дифференциального тока срабатывания «Iср (ДТЗт)». Второй участок представлен в виде наклонной линии и определяется уставками начального тормозного тока «Iт1» и коэффициентом торможения\footnote{~Коэффициент торможения это отношение приращения дифференциального тока срабатывания к приращению тормозного тока.} «Kт1». Третий участок представлен в виде наклонной линии и определяется уставками начального тормозного тока «Iт2» и коэффициентом торможения «Kт2». Уставки для всех трех фаз задаются одинаковыми. Для исключения возможности излишнего срабатывания при насыщении измерительных ТТ на значениях тормозного тока выше «Iт2» в ДТЗт предусмотрена блокировка при внешних КЗ.
\item
Для каждой из фаз непрерывно выполняется расчет дифференциального и тормозного токов. Основной рабочей величиной ДЗТ является дифференциальный ток. В общем виде дифференциальный ток (выраженный в о.е.) равен модулю геометрической суммы токов основных гармоник плеч, определенных с учетом коэффициентов выравниваний и компенсации токов нулевой последовательности (в случае необходимости) по формуле:

% ФОРМУЛА №3
\begin{equation*}
%\label{eq:formula3}
I_{\textsl{диф.ph}} = \left | \sum_{n=1}^{m} \vec I_{ph,n*} \right |,
\end{equation*}
\begin{eqexpl}[25mm]
\item{$ph$}контролируемая фаза (A,B,C);
\item{$\vec I_{ph,n*}$}вектор тока основной гармоники $n$-го плеча, определенного с учетом коэффициентов выравниваний, о.е.;
\item{$m$}число сторон (до трёх).
\end{eqexpl}

\item
Дифференциальные токи (выраженные в о.е.) отдельных фаз определяются (в случае, если используется три стороны защищаемого объекта) по следующим формулам:  
\begin{equation*}
%\label{eq:formula4}
\begin{split}
I_{\textsl{диф.A}} = \left | \vec I_{A,1*}+\vec I_{A,2*}+\vec I_{A,3*} \right |,\\
I_{\textsl{диф.B}} = \left | \vec I_{B,1*}+\vec I_{B,2*}+\vec I_{B,3*} \right |,\\
I_{\textsl{диф.C}} = \left | \vec I_{C,1*}+\vec I_{C,2*}+\vec I_{C,3*} \right |.
\end{split}
\end{equation*}

\item
В общем виде тормозной ток ДТЗт (выраженный в о.е.) равен арифметической полусумме входных токов плеч, определенных с учетом коэффициентов выравниваний по формуле:
\begin{equation*}
%\label{eq:formula5}
I_{\textsl{торм.ph}} = \frac{1}{2} \sum_{n=1}^{m}\left |\vec I_{ph,n*} \right |,
\end{equation*}
\begin{eqexpl}[25mm]
\item{$ph$}контролируемая фаза (A,B,C);
\item{$\vec I_{ph,n*}$}вектор тока основной гармоники n-го плеча, определенного с учетом коэффициентов выравниваний, о.е.;
\item{$m$}число сторон (до трёх).
\end{eqexpl}

\item
Тормозные токи отдельных фаз определяются (в случае, если используется три стороны защищаемого объекта) по следующим формулам:  
\begin{equation*}
%\label{eq:formula4}
\begin{split}
I_{\textsl{торм.A}} = \frac{1}{2}(\left |\vec I_{A,1*}|+|\vec I_{A,2*}|+|\vec I_{A,3*} \right |),\\
I_{\textsl{торм.B}} = \frac{1}{2}(\left |\vec I_{B,1*}|+|\vec I_{B,2*}|+|\vec I_{B,3*} \right |),\\
I_{\textsl{торм.C}} = \frac{1}{2}(\left |\vec I_{C,1*}|+|\vec I_{C,2*}|+|\vec I_{C,3*} \right |).
\end{split}
\end{equation*}

\item
В случае попадания расчетных величин в зону срабатывания на время большее времени уставки «Тср» формируется сигнал срабатывания ДТЗт.

\item
Ввод ДТЗт в работу на этапе параметрирования устройства осуществляется программным переключателем SGF1 «Ввод\_функции». Информация о введенной функции отображается выходным сигналом «ДТЗт: Ввод».

\item
ДТЗт может быть выведена из работы оперативно путем активации сигнала «ОВ ДТЗт», что характеризуется активным состоянием сигнала «ДТЗт: Оперативный вывод» на выходе функции.

\item
При активном состоянии входного сигнала «ДТЗт~на~сигн» действие отключающей логики переводится на сигнал.

\item
Программным переключателем SGF2 «Реж\_блок» возможен выбор режима блокировки ДТЗт: «Без блокировки», «Блокировка по второй гармонике», «Блокировка по пятой гармонике», «Блокировка по второй и пятой гармоникам».

\item
В алгоритме ДТЗт предусмотрена работа с дополнительным подтверждением от функции блокировки от внешних КЗ (БВКЗ) путем выбора состояния «С контролем БВКЗ» программного переключателя SGF3 «Контр\_БВКЗ», описание работы БВКЗ приведено в подразделе \ref{sec:bvkz}.

\item
При выявлении неисправности в цепях ТТ, что определяется срабатыванием КЦТ по сторонам трансформатора (см. подраздел \ref{sec:kzt}) автоматически увеличивается уровень тока срабатывания горизонтального участка тормозной характеристики для снижения вероятности ложного срабатывания ДТЗт в режиме протекания токов нагрузки путем загрубления уставки ДТЗт до уровня «Iср\_загр» (см. рисунок \ref{dzt:char}).

\item
Алгоритм работы ДТЗт выполнен в соответствии с рисунком \ref{dzt:img1}. В таблице \ref{dzt:tbl1} приведены параметры, необходимые для настройки ДТЗт.

\vspace{1mm}
\begin{figure}[H]
\includegraphics[width=1\textwidth,height=1\textheight,keepaspectratio]{\fbpath/17. ДЗТ 35/_latex/img/dtz.pdf}
\captionof{figure}{Функциональная схема логики ДТЗт}
\label{dzt:img1}
\end{figure}
\small
\begin{longtable}{|>{\centering\arraybackslash}m{5.3cm}|>{\centering\arraybackslash}m{3.3cm}|>{\centering\arraybackslash}m{4.2cm}|>{\centering\arraybackslash}m{1.8cm}|>{\centering\arraybackslash}m{1cm}|}
\caption{Параметры для настройки ДТЗт\hfill\vspace{-0.5\baselineskip}}\label{dzt:tbl1}\\ 
\hline
\rowcolor{gray!30}
Параметр (Параметр на ИЧМ) & Условное обозначение на схеме & Значение/ Диапазон & Единица измерения & Шаг \\ 
\hline
\endfirsthead
%\caption{\emph{Продолжение\hfill\vspace{-0.5\baselineskip}}} \\ % Отступ обозначения таблицы от таблицы и выравнивание надписи слева
\caption*{\hspace{3pt}\emph{Продолжение таблицы \ref{dzt:tbl1}\hfill\vspace{-0.5\baselineskip}}} \\ % сделано по ГОСТ 2.105 п.6.8.7
\hline
\rowcolor{gray!30}
Параметр (Параметр на ИЧМ) & Условное обозначение на схеме & Значение/ Диапазон & Единица измерения & Шаг \\ 
%\hline
\endhead
%\hline
\endfoot
%\hline
\endlastfoot
%==+t1*PDIF1|TDIF
%===t1*PDIF1
\centering Ввод функции в работу (Ввод\_функции) & \centering SGF1 & \centering 0 = Не предусмотрено\\1 = Предусмотрено & \centering -- & \centering \arraybackslash -- \\
\hline
\centering Контроль от БВКЗ (Контр\_БВКЗ) & \centering SGF3 & \centering 0 = Без контроля БВКЗ\\1 = С контролем БВКЗ & \centering -- & \centering \arraybackslash -- \\
\hline
\centering Режим блокировки (Реж\_блок) & \centering SGF2 & \centering 1 = Без блокировки\\2 = Блокировка по 2 гармонике\\3 = Блокировка по 5 гармонике\\4 = Блокировка по 2 и 5 гармоникам & \centering -- & \centering \arraybackslash -- \\
\hline
\centering Начальный дифференциальный ток срабатывания (Iср) & \centering -- & \centering 0,20 ... 0,60 & \centering о.е. & \centering \arraybackslash 0,01 \\
\hline
\centering Начальный дифференциальный ток срабатывания в режиме загрубления (при неисправностях в токовых цепях) (Iср\_загр) & \centering -- & \centering 0,20 ... 10,00 & \centering о.е. & \centering \arraybackslash 0,01 \\
\hline
\centering Коэффициент торможения первого наклонного участка (Kт1) & \centering -- & \centering 0,20 ... 1,00 & \centering о.е. & \centering \arraybackslash 0,01 \\
\hline
\centering Начальный тормозной ток первого наклонного участка (Iт1) & \centering -- & \centering 0,60 ... 3,00 & \centering о.е. & \centering \arraybackslash 0,01 \\
\hline
\centering Коэффициент торможения второго наклонного участка (Kт2) & \centering -- & \centering 0,20 ... 1,00 & \centering о.е. & \centering \arraybackslash 0,01 \\
\hline
\centering Начальный тормозной ток второго наклонного участка (Iт2) & \centering -- & \centering 1,20 ... 10,00 & \centering о.е. & \centering \arraybackslash 0,01 \\
\hline
\centering Выдержка времени срабатывания (Tср) & \centering T1 & \centering 0,00 ... 3,00 & \centering с & \centering \arraybackslash 0,01 \\
\hline
%===t1
\end{longtable}
\normalsize

\end{enumerate}
%%%%%%%%%%%%%%%%%%%%%%%%%%%%%%%%%%%%%%%%%%%%%%%%%%%%%%%%%%%%%%%%%%%%%%  ДИФОТСЕЧКА %%%%%%%%%%%%%%%%%%%%%%%%%%%%%%%%%%%%%%%%%%%%%%%%%%

%\needspace{3\baselineskip}
%\color{uniblue}{\subsubsection{Дифференциальная токовая отсечка (ДТО)}}
%\color{black}
\item Дифференциальная токовая отсечка (ДТО)

\begin{enumerate}[label=\arabic{section}.\arabic{subsection}.\arabic{enumi}.\arabic*, labelsep=4pt, leftmargin=0em, itemindent=65pt, parsep=0pt]

\item
Дифференциальная токовая отсечка (ДТО) предназначена для быстрого и селективного отключения повреждений, сопровождающихся большими токами КЗ в зоне действия ДЗТ. Характеристика срабатывания ДТО (см. рисунок \ref{dzt:char}) представляет собой прямую линию, определяемую уставкой «Iср». ДТО работает без каких-либо блокировок и не имеет торможения. Срабатывание ДТО происходит, если дифференциальный ток превышает уставку «Iср» в течении выдержки времени «Тср». Возврат ДТО происходит, если дифференциальный ток ниже уставки «Iср» с учетом коэффициента возврата («Kвоз» равен 0,95).

\item
Ввод ДТО в работу на этапе параметрирования устройства осуществляется программным переключателем SGF1 «Ввод\_функции». Информация о введенной функции отображается выходным сигналом «ДТО: Ввод».
\item
ДТО может быть выведена из работы оперативно путем активации сигнала «ОВ ДТО», что характеризуется активным состоянием сигнала «ДТО: Оперативный вывод» на выходе функции. 
\item
При активном состоянии входного сигнала «ДТО~на~сигн» действие отключающей логики переводится на сигнал.
\item
Алгоритм работы ДТО выполнен в соответствии с рисунком \ref{dzt:img2}. В таблице \ref{dzt:tbl2} приведены параметры, необходимые для настройки ДТО.

\small
\begin{longtable}{|>{\centering\arraybackslash}m{5.3cm}|>{\centering\arraybackslash}m{3.3cm}|>{\centering\arraybackslash}m{4.2cm}|>{\centering\arraybackslash}m{1.8cm}|>{\centering\arraybackslash}m{1cm}|}
\caption{Параметры для настройки ДТО\hfill\vspace{-0.5\baselineskip}}\label{dzt:tbl2}\\ 
\hline
\rowcolor{gray!30}
Параметр (Параметр на ИЧМ) & Условное обозначение на схеме & Значение/ Диапазон & Единица измерения & Шаг \\ 
\hline
\endfirsthead
%\caption{\emph{Продолжение\hfill\vspace{-0.5\baselineskip}}} \\ % Отступ обозначения таблицы от таблицы и выравнивание надписи слева
\caption*{\hspace{3pt}\emph{Продолжение таблицы \ref{dzt:tbl2}\hfill\vspace{-0.5\baselineskip}}} \\ % сделано по ГОСТ 2.105 п.6.8.7
\hline
\rowcolor{gray!30}
Параметр (Параметр на ИЧМ) & Условное обозначение на схеме & Значение/ Диапазон & Единица измерения & Шаг \\ 
%\hline
\endhead
%\hline
\endfoot
%\hline
\endlastfoot
%==+t1*PDIF2|TDIF
%===t1*PDIF2
\centering Ввод функции в работу (Ввод\_функции) & \centering SGF1 & \centering 0 = Не предусмотрено\\1 = Предусмотрено & \centering -- & \centering \arraybackslash -- \\
\hline
\centering Дифференциальный ток срабатывания (Iср) & \centering I> & \centering 3,0 ... 20,0 & \centering о.е. & \centering \arraybackslash 0,1 \\
\hline
\centering Выдержка времени срабатывания (Tср) & \centering T1 & \centering 0,00 ... 3,00 & \centering с & \centering \arraybackslash 0,01 \\
\hline
%===t1
\end{longtable}
\normalsize

\vspace{1mm}
\begin{figure}[H]
\includegraphics[width=\textwidth,height=\textheight,keepaspectratio]{\fbpath/17. ДЗТ 35/_latex/img/dto.pdf}
\captionof{figure}{Функциональная схема логики ДТО}
\label{dzt:img2}
\end{figure}


\end{enumerate}
%%%%%%%%%%%%%%%%%%%%%%%%%%%%%%%%%%%%%%%%%%%%%%%%%%%%%%%%%%%%%%%%%%%%%%  ДЕТЕКТОР 2 ГАРМОНИКИ %%%%%%%%%%%%%%%%%%%%%%%%%%%%%%%%%%%%%%%%%%%%%%%%%%

%\needspace{3\baselineskip}
%\color{uniblue}{\subsubsection{Детектор второй гармоники (Д2Г)}}
%\color{black}

\item Детектор второй гармоники (Д2Г)

\begin{enumerate}[label=\arabic{section}.\arabic{subsection}.\arabic{enumi}.\arabic*, labelsep=4pt, leftmargin=0em, itemindent=65pt, parsep=0pt]

\item
Детектор второй гармоники (Д2Г) предусмотрен для отстройки ДЗТ от режимов бросков токов намагничивания, обусловленными остаточным насыщением магнитопровода при включении силового трансформатора. Алгоритм Д2Г контролирует уровень второй гармоники в дифференциальном токе и реагирует на отношение модуля второй гармоники дифференциального тока к модулю основной гармоники. Если данное соотношение превышает уставку «Ih2/Ih1», то происходит блокировка срабатывания ДТЗт, которая сохраняется до тех пор, пока соотношение не опустится ниже заданной уставки с учетом коэффициента возврата после набора выдержки времени на возврат «Твоз». Срабатывание ДТЗт блокируется пофазно при появлении соответствующего блокирующего сигнала. 
\item
Ввод Д2Г в работу осуществляется автоматически при выполнении следующих условий:
\begin{itemize}
\item ввод ДТЗт в работу; 
\item\sloppy посредством программного переключателя «Режим блокировки» ДТЗт выбран один из режимов: «Блокировка по 2 гармонике» или «Блокировка по 2 и 5 гармоникам».
\end{itemize}
\item
Дополнительно в логике Д2Г должен быть реализован режим перекрестной блокировки (ввод производится уставкой «Перекр\_блок»). Принцип работы данного режима заключается в следующем:
\begin{itemize}
\item при обнаружении условий срабатывания Д2Г в одной из трех фаз должно произойти блокирование срабатывания ДТЗт по всем трем фазам на время, задаваемое уставкой «Tперек\_блок»;
\item при пропадании условий срабатывания Д2Г от смежных фаз или по истечению времени «Tперек\_блок» должно осуществляться блокирование срабатывания ДТЗт пофазно.
\end{itemize}
\item
Алгоритм работы Д2Г выполнен в соответствии с рисунком \ref{dzt:imgd2g}. В таблице \ref{dzt:tbld2g} приведены параметры, необходимые для настройки Д2Г. 

\vspace{1mm}
\begin{figure}[H]
\includegraphics[width=\textwidth,height=\textheight,keepaspectratio]{\fbpath/17. ДЗТ 35/_latex/img/d2g.pdf}
\captionof{figure}{Функциональная схема логики Д2Г}
\label{dzt:imgd2g}
\end{figure}

\small
\begin{longtable}{|>{\centering\arraybackslash}m{5.3cm}|>{\centering\arraybackslash}m{3.3cm}|>{\centering\arraybackslash}m{4.2cm}|>{\centering\arraybackslash}m{1.8cm}|>{\centering\arraybackslash}m{1cm}|}
\caption{Параметры для настройки Д2Г\hfill\vspace{-0.5\baselineskip}}\label{dzt:tbld2g}\\ 
\hline
\rowcolor{gray!30}
Параметр (Параметр на ИЧМ) & Условное обозначение на схеме & Значение/ Диапазон & Единица измерения & Шаг \\ 
\hline
\endfirsthead
%\caption{\emph{Продолжение\hfill\vspace{-0.5\baselineskip}}} \\ % Отступ обозначения таблицы от таблицы и выравнивание надписи слева
\caption*{\hspace{3pt}\emph{Продолжение таблицы \ref{dzt:tbld2g}\hfill\vspace{-0.5\baselineskip}}} \\ % сделано по ГОСТ 2.105 п.6.8.7
\hline
\rowcolor{gray!30}
Параметр (Параметр на ИЧМ) & Условное обозначение на схеме & Значение/ Диапазон & Единица измерения & Шаг \\ 
%\hline
\endhead
%\hline
\endfoot
%\hline
\endlastfoot
%==+t1*HF2PHAR1|TDIF
%===t1*HF2PHAR1
\centering Коэффициент отношения второй гармоники к первой гармонике для выполнения пусковых условий (Ih2/Ih1) & \centering I2г/I1г> & \centering 0,10 ... 0,50 & \centering -- & \centering \arraybackslash 0,01 \\
\hline
\centering Режим перекрестной блокировки (Перекр\_блок) & \centering SGF1 & \centering 0 = Не предусмотрено\\1 = Предусмотрено & \centering -- & \centering \arraybackslash -- \\
\hline
\centering Выдержка времени ввода перекрестной блокировки (Tперек\_блок) & \centering T2 & \centering 0,06 ... 4,00 & \centering с & \centering \arraybackslash 0,01 \\
\hline
\centering Выдержка времени на возврат (Tвоз) & \centering T1 & \centering 0,000 ... 0,100 & \centering с & \centering \arraybackslash 0,001 \\
\hline
%===t1
\end{longtable}
\normalsize

\end{enumerate}
%%%%%%%%%%%%%%%%%%%%%%%%%%%%%%%%%%%%%%%%%%%%%%%%%%%%%%%%%%%%%%%%%%%%%%  ДЕТЕКТОР 5 ГАРМОНИКИ %%%%%%%%%%%%%%%%%%%%%%%%%%%%%%%%%%%%%%%%%%%%%%%%%%

%\needspace{3\baselineskip}
%\color{uniblue}{\subsubsection{Детектор пятой гармоники (Д5Г)}}
%\color{black}

\item Детектор пятой гармоники (Д5Г)

\begin{enumerate}[label=\arabic{section}.\arabic{subsection}.\arabic{enumi}.\arabic*, labelsep=4pt, leftmargin=0em, itemindent=65pt, parsep=0pt]

\item
Детектор пятой гармоники (Д5Г) предусмотрен для отстройки ДЗТ от режимов бросков токов намагничивания, обусловленными остаточным насыщением магнитопровода при включении силового трансформатора. Алгоритм Д5Г контролирует уровень пятой гармоники в дифференциальном токе и реагирует на отношение модуля второй гармоники дифференциального тока к модулю основной гармоники. Если данное соотношение превышает уставку «Ih5/Ih1», то происходит блокировка срабатывания ДТЗт, которая сохраняется до тех пор, пока соотношение не опустится ниже заданной уставки с учетом коэффициента возврата после набора выдержки времени на возврат «Твоз». Срабатывание ДТЗт блокируется пофазно при появлении соответствующего блокирующего сигнала. 
\item
Ввод Д5Г в работу осуществляется автоматически при выполнении следующих условий:
\begin{itemize}
\item ввод ДТЗт в работу; 
\item\sloppy посредством программного переключателя «Режим блокировки» ДТЗт выбран один из режимов: «Блокировка по 5 гармонике» или «Блокировка по 2 и 5 гармоникам».
\end{itemize}
\item
Дополнительно, в логике Д5Г должен быть реализован режим перекрестной блокировки (ввод производится уставкой «Перекр\_блок»). Принцип работы данного режима заключается в следующем:
\begin{itemize}
\item при обнаружении условий срабатывания Д5Г в одной из трех фаз должно произойти блокирование срабатывания ДТЗт по всем трем фазам на время, задаваемое уставкой «Tперек\_блок»;
\item при пропадании условий срабатывания Д5Г от смежных фаз или по истечению времени «Tперек\_блок» должно осуществляться блокирование срабатывания ДТЗт пофазно.
\end{itemize}

\item
Алгоритм работы Д5Г выполнен в соответствии с рисунком \ref{dzt:imgd5g}. В таблице \ref{dzt:tbld5g} приведены параметры, необходимые для настройки Д5Г.

\vspace{1mm}
\begin{figure}[H]
\includegraphics[width=\textwidth,height=\textheight,keepaspectratio]{\fbpath/17. ДЗТ 35/_latex/img/d5g.pdf}
\captionof{figure}{Функциональная схема логики Д5Г}
\label{dzt:imgd5g}
\end{figure}

\small
\begin{longtable}{|>{\centering\arraybackslash}m{5.3cm}|>{\centering\arraybackslash}m{3.3cm}|>{\centering\arraybackslash}m{4.2cm}|>{\centering\arraybackslash}m{1.8cm}|>{\centering\arraybackslash}m{1cm}|}
\caption{Параметры для настройки Д5Г\hfill\vspace{-0.5\baselineskip}}\label{dzt:tbld5g}\\ 
\hline
\rowcolor{gray!30}
Параметр (Параметр на ИЧМ) & Условное обозначение на схеме & Значение/ Диапазон & Единица измерения & Шаг \\ 
\hline
\endfirsthead
%\caption{\emph{Продолжение\hfill\vspace{-0.5\baselineskip}}} \\ % Отступ обозначения таблицы от таблицы и выравнивание надписи слева
\caption*{\hspace{3pt}\emph{Продолжение таблицы \ref{dzt:tbld5g}\hfill\vspace{-0.5\baselineskip}}} \\ % сделано по ГОСТ 2.105 п.6.8.7
\hline
\rowcolor{gray!30}
Параметр (Параметр на ИЧМ) & Условное обозначение на схеме & Значение/ Диапазон & Единица измерения & Шаг \\ 
%\hline
\endhead
%\hline
\endfoot
%\hline
\endlastfoot
%==+t1*HF5PHAR1|TDIF
%===t1*HF5PHAR1
\centering Коэффициент отношения пятой гармоники к первой гармонике для выполнения пусковых условий (Ih5/Ih1) & \centering I5г/I1г> & \centering 0,05 ... 0,40 & \centering -- & \centering \arraybackslash 0,01 \\
\hline
\centering Режим перекрестной блокировки (Перекр\_блок) & \centering SGF1 & \centering 0 = Не предусмотрено\\1 = Предусмотрено & \centering -- & \centering \arraybackslash -- \\
\hline
\centering Выдержка времени ввода перекрестной блокировки (Tперек\_блок) & \centering T2 & \centering 0,06 ... 4,00 & \centering с & \centering \arraybackslash 0,01 \\
\hline
\centering Выдержка времени на возврат (Tвоз) & \centering T1 & \centering 0,000 ... 0,100 & \centering с & \centering \arraybackslash 0,001 \\
\hline
%===t1
\end{longtable}
\normalsize

\end{enumerate}
%%%%%%%%%%%%%%%%%%%%%%%%%%%%%%%%%%%%%%%%%%%%%% БЛОКИРОВКА ПРИ ВНЕШНИХ КЗ %%%%%%%%%%%%%%%%%%%%%%%%%%%%%%%%%%%%%%%%%%%%%%%

%\needspace{3\baselineskip}
%\color{uniblue}{\subsubsection{Блокировка при внешних КЗ (БВКЗ)}\label{sec:bvkz}}
%\color{black}

\item\label{sec:bvkz} Блокировка при внешних КЗ (БВКЗ)

\begin{enumerate}[label=\arabic{section}.\arabic{subsection}.\arabic{enumi}.\arabic*, labelsep=4pt, leftmargin=0em, itemindent=65pt, parsep=0pt]

\item
Срабатывание блокировки при внешних КЗ (БВКЗ) является дополнительным критерием стабильной работы ДЗТ при насыщении ТТ. Общий принцип действия состоит в сравнении фаз токов по сторонам защищаемого объекта. БВКЗ выполняется в пофазном исполнении.
\item
Временные диаграммы работы БВКЗ в режимах внешнего и внутреннего КЗ приведены на рисунке \ref{dzt:imgcharbvkz}, для однозначного соответствия точек измерения на функциональной схеме (рис. \ref{dzt:imgbvkz}) и на временных диаграммах (рис. \ref{dzt:imgcharbvkz}) точкам измерения присвоены цифры \textcircled{\raisebox{-0.9pt}{1}}...\textcircled{\raisebox{-0.9pt}{5}}.
На функциональной схеме (рис. \ref{dzt:imgbvkz}) и временных диаграммах (рис. \ref{dzt:imgcharbvkz}) указаны точки измерений. 

%\emph{\textls[50]{Выбор величин для сравнения}}
\item
Каждый период выполняется вычисление амплитуд токов сторон. Ток стороны с максимальной амплитудой выбирается в качестве первой величины для сравнения ($i1$). В качестве второй величины для сравнения ($i2$) используется сумма мгновенных значений остальных сторон.

%\emph{\textls[50]{Сравнение фаз}}
\item
Во время положительной полуволны тока каждой из сравниваемых величин устанавливаются сигналы 1 на определенных элементах схемы. Далее эти сигналы суммируются. При внешних КЗ в сумме сигналов практически отсутствуют паузы (паузы обусловлены пороговыми значениями элементов сравнения). При внутренних КЗ паузы по длительности приближаются к значению времени полупериода (см. рисунок \ref{dzt:imgcharbvkz}). Длительность паузы является критерием блокировки при внешних КЗ. Вторая по счету от начала повреждения длительная пауза дает разрешение на работу ДЗТ.

\vspace{1mm}
\begin{figure}[H]
\centering
\includegraphics[width=1\textwidth,height=1\textheight,keepaspectratio]{\fbpath/17. ДЗТ 35/_latex/img/charbvkz.pdf}
\captionof{figure}{Временные диаграммы работы БВКЗ}
\label{dzt:imgcharbvkz}
\end{figure}

\item
БВКЗ автоматически вводится в работу при введенной в работу ДТЗт и введенном переключателе SGF3 «Контр\_БВКЗ» в составе ДТЗт.

\item
Алгоритм работы БВКЗ выполнен в соответствии с рисунком \ref{dzt:imgbvkz}. БВКЗ не имеет настроечных параметров.

\vspace{1mm}
\begin{figure}[H]
\includegraphics[width=\textwidth,height=\textheight,keepaspectratio]{\fbpath/17. ДЗТ 35/_latex/img/bvkz.pdf}
\captionof{figure}{Функциональная схема логики БВКЗ}
\label{dzt:imgbvkz}
\end{figure}

\end{enumerate}
%%%%%%%%%%%%%%%%%%%%%%%%%%%%%%%%%%%%%%%%%%%%%%%%%%%%%%%%%%%%%%%%% КОНТРОЛЬ ИСПРАВНОСТИ ТОКОВЫХ ЦЕПЕЙ ПО НЕБАЛАНСУ %%%%%%%%%%%%%%%%%%%%%%%%%%%%%%%%%%%%%%%%%%%%%%%%

%\needspace{3\baselineskip}
%\color{uniblue}{\subsubsection{Контроль исправности токовых цепей по величине небаланса (КЦТнеб)}}
%\color{black}

\item Контроль исправности токовых цепей по величине небаланса (КЦТнеб)

\begin{enumerate}[label=\arabic{section}.\arabic{subsection}.\arabic{enumi}.\arabic*, labelsep=4pt, leftmargin=0em, itemindent=65pt, parsep=0pt]

\item
Контроль исправности токовых цепей по величине небаланса (КЦТнеб) в дифференциальных цепях является дополнительным. Принцип действия этой функции основан на измерении тока небаланса в дифференциальных цепях: если дифференциальный ток превышает уставку тока небаланса «Iср» в течение выдержки времени «Tср», то выдается сигнал о неисправности в сигнализацию.
\item
Ввод функции в работу на этапе параметрирования устройства осуществляется программным переключателем SGF1 «Ввод\_функции». Информация о введенной функции отображается выходным сигналом «КЦТнеб: Ввод».
\item
КЦТнеб может быть выведен из работы оперативно путем активации сигнала «ОВ КЦТнеб», что характеризуется активным состоянием сигнала «КЦТнеб: Оперативный вывод» на выходе функции. 
\item
Алгоритм работы КЦТнеб выполнен в соответствии с рисунком \ref{dzt:imgkzt}. В таблице \ref{dzt:tblkzt} приведены параметры, необходимые для настройки КЦТнеб.

\vspace{1mm}
\begin{figure}[H]
\includegraphics[width=\textwidth,height=\textheight,keepaspectratio]{\fbpath/17. ДЗТ 35/_latex/img/kzt.pdf}
\captionof{figure}{Функциональная схема логики КЦТнеб}
\label{dzt:imgkzt}
\end{figure}
\small
\begin{longtable}{|>{\centering\arraybackslash}m{5.3cm}|>{\centering\arraybackslash}m{3.3cm}|>{\centering\arraybackslash}m{4.2cm}|>{\centering\arraybackslash}m{1.8cm}|>{\centering\arraybackslash}m{1cm}|}
\caption{Параметры для настройки КЦТнеб\hfill\vspace{-0.5\baselineskip}}\label{dzt:tblkzt}\\ 
\hline
\rowcolor{gray!30}
Параметр (Параметр на ИЧМ) & Условное обозначение на схеме & Значение/ Диапазон & Единица измерения & Шаг \\ 
\hline
\endfirsthead
%\caption{\emph{Продолжение\hfill\vspace{-0.5\baselineskip}}} \\ % Отступ обозначения таблицы от таблицы и выравнивание надписи слева
\caption*{\hspace{3pt}\emph{Продолжение таблицы \ref{dzt:tblkzt}\hfill\vspace{-0.5\baselineskip}}} \\ % сделано по ГОСТ 2.105 п.6.8.7
\hline
\rowcolor{gray!30}
Параметр (Параметр на ИЧМ) & Условное обозначение на схеме & Значение/ Диапазон & Единица измерения & Шаг \\ 
%\hline
\endhead
%\hline
\endfoot
%\hline
\endlastfoot
%==+t1*RCTR1|TDIF
%===t1*RCTR1
\centering Ввод функции в работу (Ввод\_функции) & \centering SGF1 & \centering 0 = Не предусмотрено\\1 = Предусмотрено & \centering -- & \centering \arraybackslash -- \\
\hline
\centering Ток срабатывания (Iср) & \centering Idiff> & \centering 0,04 ... 2,00 & \centering о.е. & \centering \arraybackslash 0,01 \\
\hline
\centering Выдержка времени срабатывания (Tср) & \centering T1 & \centering 0,0 ... 110,0 & \centering с & \centering \arraybackslash 0,1 \\
\hline
%===t1
\end{longtable}
\normalsize
\end{enumerate}

%%%%%%%%%%%%%%%%%%%%%%%%%%%%%%%%%%%%%%%%%%%%%%%%%%%%%%%%%%%%%%%%% ОБЩИЕ УСТАВКИ %%%%%%%%%%%%%%%%%%%%%%%%%%%%%%%%%%%%%%%%%%%%%%%%
\item Общие параметры 

\begin{enumerate}[label=\arabic{section}.\arabic{subsection}.\arabic{enumi}.\arabic*, labelsep=4pt, leftmargin=0em, itemindent=65pt, parsep=0pt]

\item
Для правильной работы ДЗТ необходимо задать общие параметры (базисная мощность трансформатора, схема соединения ТТ для каждой из сторон, выбор сторон для расчета дифференциального и тормозного тока). 
В таблице \ref{dzt:tblgen} приведены общие параметры, необходимые для настройки ДЗТ.

\small
\begin{longtable}{|>{\centering\arraybackslash}m{5.3cm}|>{\centering\arraybackslash}m{3.3cm}|>{\centering\arraybackslash}m{4.2cm}|>{\centering\arraybackslash}m{1.8cm}|>{\centering\arraybackslash}m{1cm}|}
\caption{Общие параметры для настройки ДЗТ\hfill\vspace{-0.5\baselineskip}}\label{dzt:tblgen}\\ 
\hline
\rowcolor{gray!30}
Параметр (Параметр на ИЧМ) & Условное обозначение на схеме & Значение/ Диапазон & Единица измерения & Шаг \\ 
\hline
\endfirsthead
\caption*{\hspace{3pt}\emph{Продолжение таблицы \ref{dzt:tblgen}\hfill\vspace{-0.5\baselineskip}}} \\ % сделано по ГОСТ 2.105 п.6.8.7
\hline
\rowcolor{gray!30}
Параметр (Параметр на ИЧМ) & Условное обозначение на схеме & Значение/ Диапазон & Единица измерения & Шаг \\ 
%\hline
\endhead
%\hline
\endfoot
%\hline
\endlastfoot
%==+t1*LLN0|TDIF
%===t1*GENERAL
\centering Сторона 1 (Сторона 1) & \centering -- & \centering 0 = Не используется\\1 = Используется & \centering -- & \centering \arraybackslash -- \\
\hline
\centering Сторона 2 (Сторона 2) & \centering -- & \centering 0 = Не используется\\1 = Используется & \centering -- & \centering \arraybackslash -- \\
\hline
\centering Сторона 3 (Сторона 3) & \centering -- & \centering 0 = Не используется\\1 = Используется & \centering -- & \centering \arraybackslash -- \\
\hline
\centering Схема соединения трансформаторов тока стороны 1 (KсхемСт1) & \centering -- & \centering 0 = Звезда\\1 = Треугольник & \centering -- & \centering \arraybackslash -- \\
\hline
\centering Схема соединения трансформаторов тока стороны 2 (KсхемСт2) & \centering -- & \centering 0 = Звезда\\1 = Треугольник & \centering -- & \centering \arraybackslash -- \\
\hline
\centering Схема соединения трансформаторов тока стороны 3 (KсхемСт3) & \centering -- & \centering 0 = Звезда\\1 = Треугольник & \centering -- & \centering \arraybackslash -- \\
\hline
\centering Базисная мощность (Sbase) & \centering -- & \centering 1000 ... 500000 & \centering кВ×А & \centering \arraybackslash 100 \\
\hline
%===t1
\end{longtable}
\normalsize
\end{enumerate}

%%%%%%%%%%%%%%%%%%%%%%%%%%%%%%%%%%%%%%%%%%%%%%%%%%%%%%%%%%%%%%%%% ПАРАМЕТРЫ НАСТРОЙКИ СТОРОН %%%%%%%%%%%%%%%%%%%%%%%%%%%%%%%%%%%%%%%%%%%%%%%%
\item Параметры настройки сторон 

\begin{enumerate}[label=\arabic{section}.\arabic{subsection}.\arabic{enumi}.\arabic*, labelsep=4pt, leftmargin=0em, itemindent=65pt, parsep=0pt]

\item
Для корректного приведения токов сторон силового трансформатора необходимо задать параметры каждой стороны (номинальное напряжение стороны, номинальный первичный ток ТТ стороны, номинальный вторичный ток ТТ стороны, требуется ли компенсация токов нулевой последовательности и номер векторной группы). 
В таблице \ref{dzt:tblside} приведены параметры, необходимые для корректного приведения токов сторон защищаемого трансформатора.

\small
\begin{longtable}{|>{\centering\arraybackslash}m{5.3cm}|>{\centering\arraybackslash}m{3.3cm}|>{\centering\arraybackslash}m{4.2cm}|>{\centering\arraybackslash}m{1.8cm}|>{\centering\arraybackslash}m{1cm}|}
\caption{Параметры, необходимые для корректного приведения токов сторон\hfill\vspace{-0.5\baselineskip}}\label{dzt:tblside}\\ 
\hline
\rowcolor{gray!30}
Параметр (Параметр на ИЧМ) & Условное обозначение на схеме & Значение/ Диапазон & Единица измерения & Шаг \\ 
\hline
\endfirsthead
\caption*{\hspace{3pt}\emph{Продолжение таблицы \ref{dzt:tblside}\hfill\vspace{-0.5\baselineskip}}} \\ % сделано по ГОСТ 2.105 п.6.8.7
\hline
\rowcolor{gray!30}
Параметр (Параметр на ИЧМ) & Условное обозначение на схеме & Значение/ Диапазон & Единица измерения & Шаг \\ 
%\hline
\endhead
%\hline
\endfoot
%\hline
\endlastfoot
%==+t1*RMXU1|TDIF>Сторона 1
%===t1*RMXU1>Сторона 1
\multicolumn{5}{|c|}{ Сторона 1 } \\ \hline 
\centering Номинальное напряжение стороны 1 (Uном1) & \centering -- & \centering 6,0 ... 250,0 & \centering кВ & \centering \arraybackslash 0,1 \\
\hline
\centering Номинальный первичный ток трансформатора тока стороны 1 (IпервСт1) & \centering -- & \centering 1 ... 40000 & \centering А & \centering \arraybackslash 1 \\
\hline
\centering Номинальный вторичный ток трансформатора тока стороны 1 (IвторСт1) & \centering -- & \centering 0,2 ... 5,0 & \centering А & \centering \arraybackslash 0,1 \\
\hline
\centering Коэффициент коррекции полярности трансформатора тока стороны 1 (KпСт1) & \centering -- & \centering -1 = Синфазное включение\\1 = Противофазное включение & \centering -- & \centering \arraybackslash -- \\
\hline
\centering Компенсация токов 3I0 для стороны 1 (Компенс3I0ст1) & \centering -- & \centering 0 = Без компенсации токов 3I0\\1 = С компенсацией токов 3I0 & \centering -- & \centering \arraybackslash -- \\
\hline
\centering Векторная группа стороны 1 (Nсх1) & \centering -- & \centering 0 ... 11 & \centering -- & \centering \arraybackslash 1 \\
\hline
%===t1
%==+t1*RMXU2|TDIF>Сторона 2
%===t1*RMXU2>Сторона 2
\multicolumn{5}{|c|}{ Сторона 2 } \\ \hline 
\centering Номинальное напряжение стороны 2 (Uном2) & \centering -- & \centering 6,0 ... 250,0 & \centering кВ & \centering \arraybackslash 0,1 \\
\hline
\centering Номинальный первичный ток трансформатора тока стороны 2 (IпервСт2) & \centering -- & \centering 1 ... 40000 & \centering А & \centering \arraybackslash 1 \\
\hline
\centering Номинальный вторичный ток трансформатора тока стороны 2 (IвторСт2) & \centering -- & \centering 0,2 ... 5,0 & \centering А & \centering \arraybackslash 0,1 \\
\hline
\centering Коэффициент коррекции полярности трансформатора тока стороны 2 (KпСт2) & \centering -- & \centering -1 = Синфазное включение\\1 = Противофазное включение & \centering -- & \centering \arraybackslash -- \\
\hline
\centering Компенсация токов 3I0 для стороны 2 (Компенс3I0ст2) & \centering -- & \centering 0 = Без компенсации токов 3I0\\1 = С компенсацией токов 3I0 & \centering -- & \centering \arraybackslash -- \\
\hline
\centering Векторная группа стороны 2 (Nсх2) & \centering -- & \centering 0 ... 11 & \centering -- & \centering \arraybackslash 1 \\
\hline
%===t1
%==+t1*RMXU3|TDIF>Сторона 3
%===t1*RMXU3>Сторона 3
\multicolumn{5}{|c|}{ Сторона 3 } \\ \hline 
\centering Номинальное напряжение стороны 3 (Uном3) & \centering -- & \centering 6,0 ... 250,0 & \centering кВ & \centering \arraybackslash 0,1 \\
\hline
\centering Номинальный первичный ток трансформатора тока стороны 3 (IпервСт3) & \centering -- & \centering 1 ... 40000 & \centering А & \centering \arraybackslash 1 \\
\hline
\centering Номинальный вторичный ток трансформатора тока стороны 3 (IвторСт3) & \centering -- & \centering 0,2 ... 5,0 & \centering А & \centering \arraybackslash 0,1 \\
\hline
\centering Коэффициент коррекции полярности трансформатора тока стороны3 (KпСт3) & \centering -- & \centering -1 = Синфазное включение\\1 = Противофазное включение & \centering -- & \centering \arraybackslash -- \\
\hline
\centering Компенсация токов 3I0 для стороны 3 (Компенс3I0ст3) & \centering -- & \centering 0 = Без компенсации токов 3I0\\1 = С компенсацией токов 3I0 & \centering -- & \centering \arraybackslash -- \\
\hline
\centering Векторная группа стороны 3 (Nсх3) & \centering -- & \centering 0 ... 11 & \centering -- & \centering \arraybackslash 1 \\
\hline
%===t1
\end{longtable}
\normalsize
\end{enumerate}

\end{enumerate}